% !TeX root=../main.tex
% در این فایل، عنوان پایان‌نامه، مشخصات خود و چکیده پایان‌نامه را به انگلیسی، وارد کنید.

%%%%%%%%%%%%%%%%%%%%%%%%%%%%%%%%%%%%
\latinuniversity{K. N. Toosi University of Technology}
\latincollege{...}
\latinfaculty{Faculty of Electrical Engineering}
\latindepartment{Mechatronics Group}
\latinsubject{Mechatronics Engineering}
%\latinfield{field}
\latintitle{Design and Implementation of a Simulator for Magnetically Levitated Planar Motor System}
\firstlatinsupervisor{Dr.Mahdi Aliyari Shooredeli}
%\firstlatinadvisor{First Advisor}
%\secondlatinadvisor{Second Advisor}
\latinname{Alireza}
\latinsurname{Amiri}
\latinthesisdate{Summer 2026}
\latinkeywords{Magnetically levitated planar motor, real‑time simulation, fuzzy‑PID control, kinematic singularity, control allocation, ROS~2, Gazebo.}
\en-abstract{
Magnetically levitated planar motors (MLPMs) offer contactless motion, high precision, and compatibility with vacuum and cleanroom environments, making them ideal for advanced industries such as flexible manufacturing, photolithography, and biotechnology. However, research and development in this field are severely constrained by the high cost and hardware complexity of building physical prototypes. This thesis aims to create a reliable, real‑time simulation platform for MLPMs—a virtual testbed that enables the design, implementation, and evaluation of control algorithms without the need for physical hardware.
To this end, an integrated framework was developed, combining the Gazebo Harmonic physics engine, the ROS~2 Jazzy robotic middleware, and MATLAB/Simulink. The dynamic and kinematic model of the system was derived from analytical force and torque equations and cast into the wrench–current mapping matrix $\Gamma$. Current allocation was performed using the Moore–Penrose pseudo‑inverse, with realistic hardware constraints enforced: a current saturation limit of $\pm9$~A and a slew‑rate limit of $80$~A/s. A fuzzy‑PID controller with 25 rules and Gaussian membership functions was designed to adaptively tune the proportional, integral, and derivative gains in real time. A two‑phase state machine ensured a safe transition from the ground‑contact phase to levitation, and a quintic‑polynomial trajectory planner generated smooth position, velocity, and acceleration references.
Experimental results showed that the fuzzy‑PID controller achieved a settling time of approximately $1.12$~s and an overshoot below $1.25\%$ on the horizontal axes. In direct comparison with a fixed‑gain classical PID, the fuzzy‑PID reduced overshoot by $43\%$ and steady‑state error by $59.7\%$. Singularity analysis using the condition number of $\Gamma$ identified a safe operating region around $\gamma \approx \pi/4$ with a minimum condition number of $18$, and the entire horizontal workspace remained well‑conditioned (condition number $< 34$). In Manhattan‑trajectory tests, the system tracked multi‑step planar paths with near‑zero error when a segment time of $1.5$~s was allocated.
The developed platform bridges the gap between theoretical control design and physical implementation, providing a foundation for future hardware deployment with minimal architectural changes.
}
